\clearpage
\thispagestyle{empty}
\phantomsection
\addcontentsline{toc}{section}{Abstract}
\section*{Abstract}

\vspace{0.5cm}

This project presents the development of an autonomous robotic system capable of 
detecting, classifying, and sorting red and blue spherical objects in real time. 
An OAK-D~S2 RGB camera mounted on the wrist and positioned above the workspace 
when the arm is in the scan pose captures image data, which is processed on a 
Raspberry Pi~5 through a deterministic computer vision pipeline. The system then 
directs a four-degree-of-freedom arm with a fifth motor for the gripper to 
physically sort the detected balls into designated containers.

\vspace{0.4cm}

A closed-form geometric inverse kinematics solver converts
detected ball positions into motor commands in less than one
millisecond. An empirical sag compensation model was developed
to correct for gravitational droop, but was superseded in the
deployed system by per-point surface heights from a touch
calibration procedure. An adaptive grip algorithm exploits the
Dynamixel motor's built-in current and load registers to verify
whether a ball is held, enabling the system to detect and
recover from failed picks without external force sensors.

\vspace{0.4cm}

The vision pipeline combines multi-range HSV colour segmentation, Hough Circle 
Transform, and geometric contour validation into an ensemble detection system. 
Detections from multiple methods are merged using a Union-Find structure and 
filtered through Non-Maximum Suppression. Adaptive lighting compensation via 
CLAHE ensures robust performance across ambient conditions ranging from 
300 to 700\,lux. A Support Vector Machine classifier was initially evaluated 
as a secondary colour-verification step but was disabled in the final 
implementation after it was found to misclassify dark navy-blue balls as red.

\vspace{0.4cm}

A classical, deterministic approach was chosen over learning-based methods such 
as convolutional neural networks. MobileNetV2 transfer learning was evaluated 
during early development but produced inference times of 200--500\,ms per frame 
on the target hardware, making it unsuitable for real-time operation. The 
proposed pipeline achieves 20--25\,FPS with a detection rate of 98--100\,\% 
and a false positive rate below 2\,\% under laboratory conditions, without 
requiring large labelled training datasets.

\vspace{0.4cm}

The system is fully explainable: every detection rejection has a specific 
geometric or chromatic cause traceable through the pipeline, making it 
well-suited for controlled sorting environments where predictable and 
diagnosable behaviour is required.

\vspace{0.4cm}

In the final 20-ball validation session, under the tested
laboratory conditions, the system sorted \textbf{20/20 balls
correctly (100\% observed net sorting success)} with \textbf{0
unrecovered failures}. Both colour classes were completed
successfully, with 10/10 blue balls and 10/10 red balls picked
and placed in the correct bins. One additional grip-verification
rejection occurred before successful retry: grip verification
correctly rejected an unstable grasp, and retry logic recovered
the cycle on the next scan. This event did not affect the final
net sorting result.

\vspace{1cm}

\noindent\textbf{Keywords:} Computer Vision, Robotic Sorting, HSV Segmentation, 
Real-Time Systems, Raspberry Pi.

\clearpage
